% =============================================================================
\chapter{O protocolo Morphe-TCP}
\label{cap:protocolo}
% =============================================================================

Cliente e servidor conversam por um protocolo binário próprio, sobre TCP, na
porta 5000. Ele é definido em dois arquivos que precisam concordar byte a byte:
\arq{C/morphe_protocol.h}, do lado do servidor, e
\arq{Python/morphe_protocol.py}, do lado do cliente.

\section{Uma conexão por operação}

Cada operação abre a sua própria conexão, e a conexão acaba com ela:

\begin{enumerate}
  \item o cliente conecta;
  \item envia um cabeçalho de requisição de 20 bytes e, logo em seguida, o
    \emph{payload} com os dados;
  \item o servidor responde com um cabeçalho de resposta de 20 bytes e o
    \emph{payload} do resultado (ou uma mensagem de erro);
  \item o servidor fecha a conexão.
\end{enumerate}

Não existe sessão, autenticação nem estado guardado entre conexões. Isso é o
que permite ao cliente escolher uma placa diferente a cada operação e a vários
clientes intercalarem operações na mesma placa (capítulo~\ref{cap:placas}). A
única exceção ao roteiro é a captura contínua do ADC, em que a resposta é uma
sequência de blocos enquanto a captura dura (seção~\ref{sec:adc-continuo}).

\section{Os cabeçalhos}

Todos os campos são inteiros sem sinal em ordem de rede
(\emph{big-endian}). A Figura~\ref{fig:protocolo} mostra os dois cabeçalhos.

\begin{figure}[htbp]
  \centering
  \figura{protocolo}
  \caption{Os cabeçalhos de requisição e de resposta, em bytes.}
  \label{fig:protocolo}
\end{figure}

\begin{description}
  \item[\texttt{magic}] a assinatura: \texttt{0x4D52504D} (``MRPM'') na
    requisição e \texttt{0x4D52504E} (``MRPN'') na resposta. Um servidor que
    recebe outra coisa responde \texttt{BAD\_MAGIC} sem ler mais nada.
  \item[\texttt{version}] sempre 1.
  \item[\texttt{opcode}] a operação (Tabela~\ref{tab:opcodes}); a resposta
    repete o da requisição.
  \item[\texttt{dtype}] o tipo dos números no \emph{payload}: 1 para
    \texttt{int32}, 2 para \texttt{float32}. Na prática o cliente sempre envia
    \texttt{int32} já em ponto fixo; o \texttt{float32} existe na convolução e
    no FIR por compatibilidade com o projeto original.
  \item[\texttt{flags}] 0, exceto no ADC, em que carrega a palavra de
    configuração do conversor.
  \item[\texttt{n\_x}, \texttt{n\_h}] os tamanhos. O significado muda com a
    operação: no IIR, \texttt{n\_h} é o número de seções; no ADC, o divisor do
    relógio que dá a taxa de amostragem.
  \item[\texttt{status}] 0 quando deu certo; os outros valores estão na
    Tabela~\ref{tab:status}. Com status diferente de zero, o \emph{payload} é
    uma mensagem de texto de \texttt{n\_out} bytes.
  \item[\texttt{n\_out}] quantos valores (ou bytes de mensagem) vêm no
    \emph{payload} da resposta.
  \item[\texttt{extra}] uma palavra livre para a operação: 1 no IIR se alguma
    amostra saturou; o divisor usado no ADC; 0 nas outras.
\end{description}

\section{As operações}

\begin{table}[htbp]
  \centering
  \small
  \caption{As operações do protocolo e o que vai em cada \emph{payload}. Todos
    os valores ocupam 4 bytes.}
  \label{tab:opcodes}
  \begin{tabularx}{\textwidth}{@{}clXX@{}}
    \toprule
    & \textbf{Operação} & \textbf{Requisição} & \textbf{Resposta} \\
    \midrule
    1 & \texttt{CONV} & \(n_x\) amostras de \(x\), depois \(n_h\) de \(h\),
      Q15.16; \(1 \le n_x, n_h \le 1024\)
      & \(n_x + n_h - 1\) amostras de \(y\), Q15.16 \\
    \addlinespace
    2 & \texttt{FFT} & 1024 amostras reais em Q15.8 (\(n_x = 1024\))
      & 1024 valores complexos em \texttt{float32}, parte real e imaginária
        intercaladas, já em escala real \\
    \addlinespace
    3 & \texttt{PING} & nada & texto \texttt{chave=valor}, uma por linha
      (seção~\ref{sec:ping}) \\
    \addlinespace
    4 & \texttt{FIR} & como a \texttt{CONV}: \(x\) e os coeficientes \(h\)
      & como a \texttt{CONV} \\
    \addlinespace
    5 & \texttt{IFFT} & 1024 valores complexos em Q15.8, real e imaginário
      intercalados (\(2 \times 1024\) palavras)
      & como a \texttt{FFT}, sem o fator \(1/N\) \\
    \addlinespace
    6 & \texttt{IIR} & \(n_x \le 1024\) amostras em Q15.16, depois
      \(5 \cdot n_h\) coeficientes Q15.16 na ordem
      \(b_0, b_1, b_2, a_1, a_2\) de cada seção; \(1 \le n_h \le 16\)
      & \(n_x\) amostras Q15.16; \texttt{extra} = 1 se saturou \\
    \addlinespace
    7 & \texttt{ADC} & nada; \(n_x\) = amostras (até 32\,768), \(n_h\) =
      divisor, \texttt{flags} = configuração
      & \(n_x\) códigos do conversor, \texttt{int32}; \texttt{extra} = divisor \\
    \addlinespace
    8 & \texttt{ADC\_CONTINUO} & como o \texttt{ADC}, com \(n_x\) sem limite
      (0 = até o cliente parar)
      & cabeçalho com \(n_{out} = 0\) e depois blocos
        (seção~\ref{sec:adc-continuo}) \\
    \bottomrule
  \end{tabularx}
\end{table}

Algumas observações sobre a tabela.

\paragraph{A FFT e a IFFT recusam tamanhos diferentes de 1024.} A convolução
completa com zeros o que faltar, porque acrescentar zeros ao fim de um sinal
não muda a convolução. Com um espectro isso não vale: completar \(X[k]\) com
zeros muda o sinal que ele representa. Por isso as duas exigem exatamente 1024
pontos, e sinais de outros tamanhos são tratados no cliente
(seção~\ref{sec:blocos}).

\paragraph{A FFT devolve números em ponto flutuante.} É a única operação em que
o servidor converte: ele conhece o expoente de bloco que o núcleo da FFT
calculou (seção~\ref{sec:acel-fft}) e o aplica, junto com o fator \(2^{-8}\) do
Q15.8, antes de enviar. O cliente ainda desfaz a escala que ele mesmo aplicou
na entrada (seção~\ref{sec:normalizacao}).

\paragraph{A IFFT recebe o dobro de dados.} Um espectro é complexo por
natureza, então o \emph{payload} tem duas palavras por ponto. O servidor satura
cada uma na faixa do Q15.8 antes de escrever na FPGA.

\paragraph{A palavra de configuração do ADC.} São os seis bits que o LTC2308
recebe a cada conversão, na ordem da folha de dados:
\texttt{S/D}, \texttt{O/S}, \texttt{S1}, \texttt{S0}, \texttt{UNI},
\texttt{SLP}. Os quatro primeiros escolhem o canal e se a entrada é simples
(contra o terra) ou diferencial (entre dois canais); \texttt{UNI} escolhe
unipolar (0 a 4,095\,V) ou bipolar (\(\pm 2{,}048\)\,V). O modo de
economia de energia (\texttt{SLP}) é recusado, porque a referência interna leva
200\,ms para voltar dele. Quem monta a palavra é o \arq{Python/aquisicao.py}; a
janela só oferece o canal e o modo.

\section{A descoberta: \texttt{PING}}
\label{sec:ping}

O \texttt{PING} é só um cabeçalho, e a resposta é um texto curto que descreve a
placa e o hardware que ela tem. Um exemplo, de uma placa preparada:

\begin{center}
\begin{minipage}{0.6\textwidth}
\small\ttfamily
service=morphe\\
version=1\\
hostname=de1soc-02\\
fft\_n=1024\\
conv\_n\_max=1024\\
iir\_secoes\_max=16\\
adc\_n\_max=32768\\
adc\_fs\_max=200000\\
uptime\_s=5412\\
fpga\_preparada=1
\end{minipage}
\end{center}

\noindent (a resposta completa traz mais alguns campos, como os bits
fracionários de cada formato.) O cliente usa o \texttt{PING} para achar placas
na rede e para conferir que uma placa está preparada
(\texttt{fpga\_preparada}); placas não preparadas não são escolhidas. O campo
\texttt{adc\_n\_max} diz se o ADC pode ser usado naquela placa agora (0 quando
não pode). O tempo que a resposta leva também é usado, para escolher a placa
menos ocupada (capítulo~\ref{cap:placas}).

\section{A captura contínua do ADC}
\label{sec:adc-continuo}

A captura contínua é a única operação cuja resposta não cabe numa mensagem só:
ela pode durar minutos, e o sinal é enviado enquanto é capturado. Depois do
cabeçalho de resposta (com \texttt{status} 0, \(n_{out} = 0\) e o divisor em
\texttt{extra}), vem uma sequência de blocos, cada um com o formato
\[
  \underbrace{n}_{\texttt{uint32}}\;
  \underbrace{\textit{estado}}_{\texttt{uint32}}\;
  \underbrace{c_0\, c_1 \ldots c_{n-1}}_{n \times \texttt{int32}}
\]
O \textit{estado} de cada bloco é 0 enquanto a captura segue. O último bloco
tem \(n = 0\) e um estado que diz por que ela terminou:

\begin{center}
\small
\begin{tabular}{@{}cl@{}}
  \toprule
  \textbf{Estado} & \textbf{Significado} \\
  \midrule
  1 & fim normal: chegou ao total pedido, ou o cliente pediu para parar \\
  2 & o servidor não acompanhou e amostras foram perdidas; o que veio antes vale \\
  3 & o hardware parou de produzir amostras \\
  \bottomrule
\end{tabular}
\end{center}

Os blocos são contínuos no tempo: a última amostra de um e a primeira do
seguinte estão separadas por exatamente um período de amostragem. Para parar a
captura antes do total, o cliente envia um byte qualquer pela mesma conexão
(é o que o botão \emph{Parar} faz), ou simplesmente a fecha. Como funciona do
lado do hardware está na seção~\ref{sec:adc}.

\section{Como acrescentar uma operação}
\label{sec:novo-opcode}

Uma operação nova toca sempre os mesmos lugares, e esquecer qualquer um deles
produz um erro que não aponta para a causa:

\begin{enumerate}
  \item \arq{C/morphe_protocol.h}: o número da operação
    (\texttt{MORPHE\_OP\_*}) e, no comentário do topo, o formato do
    \emph{payload};
  \item \arq{C/morphe_server.c}: a função \texttt{handle\_*}, seguindo o roteiro
    do capítulo~\ref{cap:servidor} --- a primeira linha confere a marca de FPGA
    preparada --- e o seu \texttt{case} no \texttt{serve\_connection};
  \item \arq{Python/morphe_protocol.py}: a constante \texttt{OP\_*}, a função
    \texttt{build\_*\_request} e a \texttt{decode\_*\_response};
  \item \arq{C/morphe_config.h} e \arq{Python/morphe_config.py}, juntos, se a
    operação tiver limites próprios;
  \item se a operação usa hardware novo: o Platform Designer, o \arq{hps_0.h}
    regenerado e um bitstream novo (capítulo~\ref{cap:desenvolvimento}).
\end{enumerate}

Mudar o formato de uma operação que já existe, em vez de criar outra, quebra os
clientes instalados nos PCs até que sejam atualizados. Se for inevitável, o
lugar de sinalizar é o campo \texttt{version}.
